\section{IPO：全球大模型第一股与公司治理新阶段}

上市章节不只是资本市场故事。张鹏谈到公司规模门槛：几十人、两百人、五百人、上市公司，对应不同管理问题。几十人时要先挣钱、先建立团队信心；一两百人时开始分工，沟通对齐成本上升；几百人以后出现层级，信息传导变长；上市以后合规和治理要求更高。

\begin{figure}[H]
\centering
\includegraphics[width=0.96\textwidth]{company-stage-thresholds.png}
\caption{创业公司阶段门槛：50 人、200 人、500 人与上市公司治理。}
\end{figure}

\begin{knowledgebox}{读图：数字背后是组织阶段}
图中数字不是精确阈值，而是组织复杂度跃迁。团队越大，越不能靠创始人个人视野覆盖所有事，必须依赖机制、委员会、流程、合规和中层管理。
\end{knowledgebox}

\subsection{学院派创业的优势与惯性}

主持人问“学院派”是否准确。张鹏承认挺对，因为团队从学院里出来。学院派的优势是重视技术、研究和长期问题；惯性则是可能忽视商业化。但智谱早期在实验室里就接触客户、挣钱、做系统交付，因此相对更早学到商业化语言。

\begin{importantbox}{学院派创业的关键修正}
学院派的技术理想主义必须补上商业化能力。研究、工程、客户、报价、成本和交付，都要进入同一张表。
\end{importantbox}

\subsection{本章小结}

IPO 是智谱的新阶段。它不仅意味着融资和知名度，也意味着治理、合规、信息传导和商业化纪律的复杂度提高。

